\begin{center}
  {\fontsize{14}{18} \textbf{Abstract}}
\end{center}

Social media platforms contain vast amounts of public opinion data, yet extracting structured sentiment insights from them remains a challenge. In this thesis, I present Reddit Sentiment Explorer, a query-driven web platform that enables users to analyze sentiment across Reddit communities for arbitrary search terms. The system collects posts and comments from Reddit via the Arctic Shift archive, filters them by language, matches them against user-supplied search terms, and classifies each matched document using a large language model through the OpenAI API. Classification results use a five-point sentiment scale ranging from very negative to very positive, accompanied by confidence scores, rationale, and evidence phrases extracted from the source text. The platform aggregates these results into interactive visualizations including sentiment timelines, distribution charts, subreddit breakdowns, and heatmaps, all presented through a Plotly Dash dashboard. The backend is built with FastAPI and SQLAlchemy on PostgreSQL, with an asynchronous worker that processes query runs independently. The entire system is containerized with Docker Compose for straightforward deployment. I evaluate the platform through a system demonstration and discuss the strengths and limitations of using large language models for sentiment classification of informal, multilingual social media text. The result is a modular, extensible tool that bridges the gap between ad-hoc sentiment analysis scripts and commercial social listening platforms.

\vspace{1em}

\textbf{Keywords:} sentiment analysis, Reddit, large language models, natural language processing, web application, Python, FastAPI, Plotly Dash

\vspace{2em}

\begin{center}
  {\fontsize{14}{16} \textbf{Összefoglaló}}
\end{center}

\begin{sloppypar}
A közösségi média platformok hatalmas mennyiségű közvélemény-adatot tartalmaznak, de a strukturált érzelemfelismerési betekintések kinyerése továbbra is kihívást jelent. Ebben a dolgozatban bemutatom a Reddit Sentiment Explorer-t, egy lekérdezés-alapú webes platformot, amely lehetővé teszi a felhasználók számára, hogy tetszőleges keresési kifejezések alapján elemezzék az érzelmi tónust a Reddit közösségeken keresztül. A rendszer az Arctic Shift archívumon keresztül gyűjti a bejegyzéseket és hozzászólásokat, nyelv szerint szűri azokat, egyezteti a felhasználó által megadott keresési kifejezésekkel, majd minden egyeztetett dokumentumot egy nagy nyelvi modell segítségével osztályoz az OpenAI API-n keresztül. Az osztályozás egy ötfokú érzelmi skálát használ, amelyhez konfidencia-pontszámok, indoklás és a forrásszövegből kiemelt bizonyíték-kifejezések tartoznak. A platform interaktív vizualizációkba összesíti az eredményeket, beleértve az érzelmi idővonalakat, eloszlási diagramokat, subreddit-bontásokat és hőtérképeket, amelyeket egy Plotly Dash vezérlőpulton mutat be. Az eredmény egy moduláris, bővíthető eszköz, amely áthidalja a szakadékot az ad-hoc érzelemelemzési szkriptek és a kereskedelmi közösségi figyelő platformok között.
\end{sloppypar}

\vspace{2em}

\begin{center}
  {\fontsize{14}{16} \textbf{Rezumat}}
\end{center}

Platformele de social media conțin cantități vaste de date privind opinia publică, însă extragerea de informații structurate despre sentimente rămâne o provocare. În această lucrare, prezint Reddit Sentiment Explorer, o platformă web bazată pe interogări care permite utilizatorilor să analizeze sentimentele din comunitățile Reddit pentru termeni de căutare arbitrari. Sistemul colectează postări și comentarii de pe Reddit prin arhiva Arctic Shift, le filtrează după limbă, le potrivește cu termenii de căutare furnizați de utilizator și clasifică fiecare document folosind un model lingvistic de mari dimensiuni prin API-ul OpenAI. Clasificarea utilizează o scală de sentiment în cinci puncte, însoțită de scoruri de încredere, raționament și expresii de evidență extrase din textul sursă. Platforma agregă aceste rezultate în vizualizări interactive, inclusiv cronologii ale sentimentelor, diagrame de distribuție, defalcări pe subreddit-uri și hărți de căldură, toate prezentate printr-un tablou de bord Plotly Dash. Rezultatul este un instrument modular și extensibil care acoperă distanța dintre scripturile ad-hoc de analiză a sentimentelor și platformele comerciale de monitorizare a rețelelor sociale.
